% Part: proof-theory
% Chapter: sequent-calculus
% Section: translations

\documentclass[../../../include/open-logic-section]{subfiles}

\begin{document}

\olfileid{pt}{seq}{tr}

\olsection{Translating between \Log{G1c} and \Log{G3c}}

We mentioned that both \Log{G1c} and \Log{G3c} are sound and complete
for classical logic, i.e., they each prove all and only the sequents
true in every !!{structure}. In particular, they prove the same
sequents. We are now in a position to prove this fact without the
detour through the soundness and completeness theorems, but using
purely proof-theoretically. Such a proof provides
\emph{transformations} of !!{proof}s in \Log{G1c} into !!{proof}s in
\Log{G3c} and vice versa.

\begin{prop}
If $\Log{G1c} \Proves S$ then $\Log{G3c} \Proves S$.
\end{prop}

\begin{proof}
  We show that if $\pi$ is !!a{proof} in \Log{G1c}, then there is
  !!a{proof}~$\pi'$ of~$S$ in\Log{G3c}. We do this by induction on $n
  = \pheight{\pi}$.

  If $n = 0$, then $\pi$ is an axiom. The axiom $\lfalse \Sequent
  \quad$ is also an axiom in~\Log{G3c} (take the context $\Gamma,
  \Delta$ to be empty). \Log{G3c} allows axioms $!A \Sequent !A$ with
  arbitrary, not just atomic~$!A$. But we proved in
  \olref[ex]{prop:G3c-axioms} that all such sequents are !!{provable}
  in~\Log{G3c}.
  
  If $n>0$, $\pi$ ends in a rule~$R$. The inductive hypothesis
  guarantees that there are !!{proof}s of the premises of that rule
  in~\Log{G3c}, i.e., that the immediate sub-!!{proof}s have
  translations in~\Log{G3c}. If the rule~$R$ is also a rule in
  \Log{G3c}, we apply it to the translated !!{proof}s of the premises.
  For those rules that are no also rules of~\Log{G3c}, we have already
  shown that they are at least admissible. So, if rule~$R$ was
  $\LeftR{\land}$ or $\RightR{\lor}$, apply
  \olref[adm]{prop:lor-G3-adm}. If it was a weakening rule, apply
  \olref[adm]{prop:weak-G3c-adm}. If it was a contraction rule, apply
  \olref[inv]{prop:G3c-cont-adm}.
\end{proof}

\begin{prop}
If $\Log{G3c} \Proves S$ then $\Log{G1c} \Proves S$.
\end{prop}

\begin{proof}
  We again proceed by induction on the !!{height} of !!a{proof}~$\pi$
  of $S$. Every axiom of \Log{G3c} can be !!{prove}d in~\Log{G1c} from
  an axiom plus several applications of the \LeftR{\Weakening} and
  \RightR{\Weakening} rules:
  \[
  \Axiom$!C \fCenter !C$
  \RightLabel{\LeftR{\Weakening}}
  \Deduce$!C, \Gamma \fCenter !C$
  \RightLabel{\RightR{\Weakening}}
  \Deduce$!C, \Gamma \fCenter \Delta, !C$
  \DisplayProof
\qquad
  \Axiom$\lfalse \fCenter$
  \RightLabel{\LeftR{\Weakening}}
  \Deduce$\lfalse, \Gamma \fCenter$
  \RightLabel{\RightR{\Weakening}}
  \Deduce$\lfalse, \Gamma \fCenter \Delta$
  \DisplayProof
  \]
  If $\pi$ ends in a rule~$R$, the inductive hypothesis yields
  !!{proof}s in \Log{G1c} of the premises, and we can apply the same
  rule~$R$ to those !!{proof}s. The exceptions are the rules
  \LeftR{\land} and \RightR{\lor}, which are derivable in \Log{G1c}:
  \[
  \Axiom$!A, !B, \Gamma \fCenter \Delta$
  \RightLabel{\LeftR{\land}}
  \UnaryInf$!A \land !B, !B, \Gamma \fCenter \Delta$
  \RightLabel{\LeftR{\land}}
  \UnaryInf$!A \land !B, !A \land !B, \Gamma \fCenter \Delta$
  \RightLabel{\LeftR{\Contraction}}
  \UnaryInf$!A \land !B, \Gamma \fCenter \Delta$
    \DisplayProof
\qquad
  \Axiom$\Gamma \fCenter \Delta, !A, !B$
  \RightLabel{\RightR{\lor}}
  \UnaryInf$\Gamma \fCenter \Delta, !A \lor !B, !B$
  \RightLabel{\RightR{\lor}}
  \UnaryInf$\Gamma \fCenter \Delta, , !A \lor !B, !A \lor !B$
  \RightLabel{\RightR{\Contraction}}
  \UnaryInf$\Gamma \fCenter \Delta, !A \lor !B$
    \DisplayProof
  \]
\end{proof}

\end{document}